\part{Startup, Performance \& Diagnostics}
\label{part:startup}

A freshly installed Emacs with a hundred packages ought to start in
under a second, yet many configurations take five or ten.  The
difference is almost never the packages themselves; it is the order
in which things load, the garbage-collection strategy during init,
and a handful of frame-parameter and file-handler settings that
compound silently.

This part disassembles the startup sequence from its first
instruction to the moment the cursor blinks.
Chapter~\ref{ch:startup-lifecycle} maps the two-file split and the
hook chain; Chapter~\ref{ch:gc-tuning} explains the dual GC triggers
whose mis-configuration can cost seconds;
Chapter~\ref{ch:general-performance} surveys a dozen runtime knobs
(font-lock deferral, bidi inhibition, process-output limits) that
affect daily editing speed.
Chapter~\ref{ch:early-init-frame} and
Chapter~\ref{ch:pre-command-hook} address the two places where timing
is critical: before the first frame and before the first keypress.
The part closes with the profiling toolkit
(Chapter~\ref{ch:profiling}) and a step-by-step diagnostic playbook
(Chapter~\ref{ch:diagnostic-playbook}) that ties everything together
into a repeatable debugging workflow.

Master this material and you will be able to reason about every
millisecond between \texttt{emacs} and your first keystroke, and to
recover quickly when a package update makes something slow.

\chapter{Emacs startup lifecycle}
\label{ch:startup-lifecycle}
\minitoc\label{emacs-startup-lifecycle}
\index{startup lifecycle}\index{early-init.el}\index{init.el}\index{package-initialize}

% ── Opening signpost ──────────────────────────────────────────────
Every Emacs session begins with a carefully ordered sequence of file
loads, hook invocations, and frame creation steps.  This chapter maps
that entire sequence, from the C core dump through
\texttt{early-init.el}, \texttt{init.el}, and the hook chain that
follows, so that you can place any piece of configuration at exactly
the right moment.  By the end you will understand why
\texttt{early-init.el} exists, what each startup hook controls, and
how a file-open request cascades through dozens of internal functions
before a fontified buffer appears on screen.  The GC and
frame-parameter optimisations that depend on this sequence are
covered in Chapters~\ref{ch:gc-tuning}
and~\ref{ch:early-init-frame}.

\section{The two-file split: \texttt{early-init.el} vs \texttt{init.el}}\label{the-two-file-split-early-init.el-vs-init.el}

Emacs loads configuration in two phases. Understanding this split is the single most important thing for startup performance.

\subsection{Phase 1: \texttt{early-init.el} (Emacs 27+)}\label{phase-1-early-init.el-emacs-27}

Loaded \textbf{before} the package system initialises and \textbf{before} the first graphical frame is created. At this point:

\begin{itemize}
\item \texttt{package.el} has \textbf{not} called \texttt{package-initialize}.
\item No frame exists; \texttt{(display-graphic-p)} returns \texttt{nil}.
\item No X/Wayland/NS resources have been read.
\item \texttt{file-name-handler-alist} is active (regex matching on every file operation).
\end{itemize}

This is the right place for:

\begin{longtable}[]{@{}ll@{}}
\toprule
Setting & Why here? \\
\midrule
\texttt{package-enable-at-startup nil} & Prevent \texttt{package.el} from scanning \texttt{elpa/} \\
\texttt{gc-cons-threshold} suppression & Must be set before any lisp loads \\
\texttt{file-name-handler-alist nil} & Removes regex overhead from every \texttt{load} \\
Frame parameters (tool-bar, etc.) & Prevents frame redraw flash \\
\texttt{native-comp-jit-compilation} & Must be set before first native-comp trigger \\
\texttt{org-fold-core-style} & Must be set before org loads \\
\bottomrule
\end{longtable}

\subsection{Phase 2: \texttt{init.el}}\label{phase-2-init.el}

Loaded after \texttt{early-init.el} and after \texttt{package-initialize} (if it wasn't suppressed). This is where \texttt{use-package} declarations, keybindings, mode hooks, and everything else goes.

\subsection{The \texttt{straight.el} wrinkle}\label{the-straight.el-wrinkle}

With \texttt{straight.el} as the package manager, \texttt{package-initialize} is suppressed (\texttt{(setq package-enable-at-startup nil)} in \texttt{early-init.el}). Instead, \texttt{straight.el} bootstraps itself at the top of \texttt{init.el}, builds its autoloads cache, and registers packages. The \texttt{straight-cache-autoloads} variable controls whether autoloads are written to a single cache file (fast) or recomputed per session.

\section{Hook firing order at startup}\label{hook-firing-order-at-startup}

This is the exact sequence of events from Emacs launch to the first user interaction:

\begin{minted}{elisp}
1. C core initialises (dump, pure space, Vloadup)
2. early-init.el evaluated
3. package-initialize called (unless suppressed)
4. init.el evaluated (top-to-bottom, synchronously)
5. after-init-hook fires
6. emacs-startup-hook fires
7. [Frame is now visible]
8. window-setup-hook fires (only in GUI mode)
9. [User can interact]
10. Idle timers begin firing (run-with-idle-timer callbacks)
\end{minted}

The critical distinction:

\begin{itemize}
\item \textbf{\texttt{after-init-hook}}: Fires immediately after \texttt{init.el} finishes evaluating, but \textbf{before} the frame is fully set up. The \texttt{emacs-init-time} timer stops here. \texttt{(info "(elisp) Startup Summary")}
\item \textbf{\texttt{emacs-startup-hook}}: Fires after \texttt{after-init-hook}, after \texttt{command-line} processing is complete. This is the hook measured by \texttt{emacs-init-time}. Safe for activating modes that need a complete environment.
\item \textbf{\texttt{window-setup-hook}}: Fires after the frame and window system are fully configured. This is the \textbf{latest} hook before user interaction. Work placed here (e.g., dashboard rendering) does not appear in \texttt{emacs-init-time}.
\end{itemize}

\subsection{Practical implication}\label{practical-implication}

If you want something to run after init but contribute to \texttt{emacs-init-time}, put it on \texttt{emacs-startup-hook}. If you want it to run but \textbf{not} show up in \texttt{emacs-init-time}, put it on \texttt{window-setup-hook}. If you want it completely off the startup path, use \texttt{run-with-idle-timer}.

\section{What happens when you open a file}\label{what-happens-when-you-open-a-file}

The chain from \texttt{C-x C-f} (or \texttt{find-file}) to a stable, fontified buffer:

\begin{minted}{text}
find-file
  → find-file-noselect
    → create buffer, insert file contents
    → file-name-handler-alist matched against filename
    → after-insert-file-functions run (coding, encoding)
    → normal-mode
      → set-auto-mode
        → check file-local variables (first line, -*- mode -*- )
        → match filename against auto-mode-alist
        → match against magic-mode-alist (first bytes)
        → check interpreter-mode-alist (shebang line)
        → apply major-mode-remap-alist (tree-sitter remapping)
        → call the winning major-mode function
          → major mode runs → derived-mode hooks fire
          → after-change-major-mode-hook fires
      → hack-local-variables (apply .dir-locals.el, file-local vars)
    → find-file-hook fires
      → git-gutter, vc-refresh, etc. run here
    → after-find-file runs
  → switch-to-buffer
    → window-buffer-change-functions run
    → jit-lock fontifies the visible portion
\end{minted}

\subsection{\texttt{file-name-handler-alist} and performance}\label{file-name-handler-alist-and-performance}

Every file operation (\texttt{load}, \texttt{require}, \texttt{insert-file-contents}) matches the filename against every entry in \texttt{file-name-handler-alist}. The default list includes handlers for TRAMP (\texttt{/ssh:}), compressed files (\texttt{.gz}), and more. During startup, when hundreds of \texttt{.elc} files are loaded, this regex matching is pure waste.

The standard pattern: nil it out in \texttt{early-init.el}, restore on \texttt{emacs-startup-hook}:

\begin{minted}{elisp}
;; early-init.el
(defvar my/file-name-handler-alist file-name-handler-alist)
(setq file-name-handler-alist nil)

(add-hook 'emacs-startup-hook
          (lambda ()
            (setq file-name-handler-alist my/file-name-handler-alist)))
\end{minted}

\texttt{(info "(elisp) Magic File Names")}

\section{\texttt{normal-mode} re-entry: the hidden performance trap}\label{normal-mode-re-entry-the-hidden-performance-trap}

\texttt{normal-mode} calls \texttt{set-auto-mode}, which determines the major mode, which triggers mode hooks. If any hook triggers a file operation that itself calls \texttt{normal-mode} (e.g., byte-compiling a file on load, which opens the \texttt{.el} to read it), you get re-entrant \texttt{normal-mode} calls. This was the root cause of the \texttt{compile-angel-on-load-mode} slowdown: every \texttt{require} triggered a byte-compile check, which opened the source file, which triggered \texttt{normal-mode} → \texttt{set-auto-mode} → mode hooks → potential further \texttt{require} calls.

% ── Mental model summary ──────────────────────────────────────────
\section{Mental model}
Emacs loads configuration in two distinct phases: \texttt{early-init.el}
runs before any frames are created and before the package system
initialises, making it the only safe place for settings that must precede
those events; \texttt{init.el} handles everything else.  Once evaluation
is complete, hooks fire in a fixed order---\texttt{after-init-hook}, then
\texttt{emacs-startup-hook}, then \texttt{window-setup-hook}, and finally
idle timers, and \texttt{emacs-init-time} measures only up to the end of
\texttt{emacs-startup-hook}, so work deferred to \texttt{window-setup-hook}
or later is invisible to that metric.

Opening a file initiates its own chain of dispatch:
\texttt{find-file} calls \texttt{normal-mode}, which calls
\texttt{set-auto-mode}, which selects a major mode and fires its hooks.
If any of those hooks trigger file operations that themselves invoke
\texttt{normal-mode}, re-entrant calls can cascade, a hidden performance
trap.  Meanwhile, \texttt{file-name-handler-alist} imposes a regex match
on every \texttt{load} and \texttt{require}; nilling it during startup
eliminates that overhead entirely.
\texttt{major-mode-remap-alist} offers a lightweight way to redirect
modes (e.g.\ to tree-sitter variants) without modifying
\texttt{auto-mode-alist} at all.

On the package-management side, \texttt{straight.el} replaces the
built-in \texttt{package-initialize} with its own bootstrap, which is why
\texttt{package-enable-at-startup} must be set to \texttt{nil} in
\texttt{early-init.el}.  The overarching strategy is to place every piece
of configuration at the latest safe hook so that \texttt{emacs-init-time}
stays low and interactive startup feels snappy.
Chapter~\ref{ch:pre-command-hook} develops this idea further with the
self-removing hook pattern.
